%% System Design & Architecture Chapter
%% This file is included in main.tex as a chapter
%% The chapter heading should be defined in main.tex before including this file

\section{Introduction}

This chapter provides comprehensive technical documentation of the Debate Platform's system design
and implementation. Chapter 1 (Introduction) outlines the core research objectives, system requirements,
and technology stack. Here, the focus is on how specific design choices enable rigorous, unbiased data
collection for empirical analysis:

\begin{itemize}
    \item \textbf{Information hiding} to prevent source bias and preserve experimental validity
    \item \textbf{WebSocket architecture} to capture real-time belief updates and debate events
    \item \textbf{Between-round reflection} as a methodological tool for measuring AI persuasion
    \item \textbf{Database schema design} to record what data is collected and why it matters
\end{itemize}

For overall system architecture, refer to the appendix (Appendix A) for detailed implementation-level documentation,
including code structure, API endpoints, WebSocket event flows, and database schema definitions.

\subsection{Architecture Summary}

The Debate Platform employs a client-server architecture:
React 18.2 frontend (with WebSocket client) communicates with Node.js/Express backend (with WebSocket server),
persisting to MongoDB. Authentication is JWT-based for stateless scaling.
AI opponents are powered by OpenAI GPT-4o mini with configurable personalities (firm, balanced, open-minded)
that modulate rhetorical strategies. Real-time communication via WebSocket enables instant argument visibility
and per-round belief tracking required for bidirectional persuasion analysis.
Information hiding is enforced client-side to prevent source bias.

\subsubsection{Design Principles}

The architecture prioritizes experimental integrity and real-time data collection:

\begin{itemize}
    \item \textbf{Real-Time Responsiveness}: WebSocket for instant argument visibility and belief tracking
    \item \textbf{Data Integrity}: Complete debate records persisted with stance evolution tracked per round
    \item \textbf{Information Hiding}: Strict client-side enforcement of source anonymity (human vs. AI)
    \item \textbf{Stateless Backend}: JWT tokens enable scaling without session persistence
    \item \textbf{Asynchronous Processing}: Non-blocking AI generation preserves UI responsiveness
\end{itemize}

\subsection{Frontend Architecture}

\subsubsection{Overview}

The frontend is a React 18.2 SPA built with Vite. It communicates with the backend via REST API for
data persistence and WebSocket for real-time debate events.
UI state is managed via two primary React Context providers:
\texttt{AuthContext} (authentication and role-based access) and
\texttt{SocketContext} (WebSocket connection and room subscriptions).
Components are organized by user role: participant (debate interaction, surveys)
 and admin (monitoring, matchmaking, user management).
 The UI is designed to enforce the experimental protocol by gating argument submission,
 prompting between-round belief checks, and capturing reflection responses in a consistent format
 that maps directly to analysis variables.

\subsubsection{Data Collection Interface}

The participant UI is structured to enforce the experimental protocol and capture analysis-ready data:
\begin{itemize}
    \item Argument submission is gated by turn order to preserve round structure
    \item Between-round prompts capture belief, confidence, and perceived influence
    \item Reflection prompts elicit short paraphrases of the opponent's argument
    \item Post-debate surveys record opponent attribution and perception confidence
\end{itemize}
Admin views emphasize monitoring and data completeness (matchmaking status, debate lifecycle, and
survey completion), ensuring each debate yields a full record for analysis.

\subsection{Backend Architecture}

\subsubsection{Overview}

The backend is a Node.js/Express server implementing stateless JWT authentication for horizontal scaling.
It follows modular separation: routes (REST endpoints and WebSocket handlers),
middleware (auth, rate limiting), models (Mongoose schemas),
services (OpenAI integration, AI argument generation), and utilities (matchmaking, cleanup jobs).
Socket.io provides real-time event streaming to debate participants and admins with room-based isolation.
These services are organized to preserve data integrity: every argument, belief update, and reflection is
validated, timestamped, and persisted as part of a complete debate record suitable for empirical analysis.
Background jobs enforce experimental integrity: debate cleanup (stale entries),
guest session cleanup (session expiration), and auto-matching based on availability.

\subsubsection{Data Collection Pipeline}

The backend is organized around reliable capture of debate interactions:
routes validate and accept arguments, belief updates, and reflections; models enforce schema-level
constraints; services generate AI responses under controlled personalities; and utilities maintain
debate lifecycle integrity (matching, cleanup, turn validation). This pipeline ensures each debate
produces a complete, timestamped record suitable for downstream analysis.
\subsection{Database Design}
\subsubsection{Overview}

MongoDB's document-based model aligns with the nested structure of debate data.
Unlike relational databases, MongoDB allows flexible schemas that evolve with feature development.
Each Debate document is self-contained with all arguments, surveys, and metadata embedded,
enabling fast single-query loads of complete debate data. This design directly supports the analysis
pipeline: all variables required for persuasion measurement (arguments, per-round belief updates,
reflection text, and opponent perception) are co-located and can be exported without expensive joins.
\subsubsection{Data Models}

\paragraph{User Schema}

\textbf{User Model Fields}
\begin{verbatim}
{
  _id: ObjectId,
  username: String (unique, required. Autogenerated for guests),
  passwordHash: String (required only for non-guests),
  role: String (admin|participant, default: participant),
  isGuest: Boolean (default: false),
  createdBy: ObjectId (reference to User who created this account),
  createdAt: Date (default: now),
  lastLogin: Date (default: null)
}
\end{verbatim}

\paragraph{Debate Schema}

\textbf{Debate Model Fields}
\begin{verbatim}
{
  _id: ObjectId,
  topicId: Number,
  topicQuestion: String,
  gameMode: String (human-human|human-ai),
  player1UserId: ObjectId,
  player1Stance: String (for|against),
  player1StanceChoice: String(for|against|unsure) (if user chooses unsure, they will be randomly assigned for/against),
  player2Type: String (human|ai),
  player2UserId: ObjectId,
  player2AIModel: String,
  player2AIPrompt: String,
  player2Stance: String (for|against),
  player2StanceChoice: String(for|against|unsure) (AI/unsure human opponents will be assigned opposite of player1Stance),
  currentBelief: {
    player1: String (for|against|unsure),
    player2: String (for|against|unsure)
  }
  currentBeliefValue: {
    player1: Number (0-100) (0: for, 100: against, 50: unsure),
    player2: Number (0-100) (0: for, 100: against, 50: unsure)
  },
    status: String (waiting|active|survey_pending|completed|abandoned),
  currentRound: Number,
  maxRounds: Number,
  firstPlayer: String (for|against),
  nextTurn: String (for|against),
  aiEnabled: Boolean,
  aiResponseDelay: Number (seconds),
  aiLastResponseAt: Date,
  arguments: [{
    stance: String(for|against),
    text: String,
    round: Number,
    submittedBy: String(human|ai),
    timestamp: Date
  }],
  preDebateSurvey: {
    player1: String(firm_on_stance|convinced_of_stance|open_to_change),
    player2: String(firm_on_stance|convinced_of_stance|open_to_change),
  },
  postDebateSurvey: {
    player1: String(strengthened|slightly_strengthened|no_effect|slightly_weakened|weakened),
    player2: String(strengthened|slightly_strengthened|no_effect|slightly_weakened|weakened),
    player1StanceStrength: Number(1-7),
    player2StanceStrength: Number(1-7),
    player1StanceConfidence: Number(1-7),
    player2StanceConfidence: Number(1-7),
    player1OpponentPerception: String(human|ai|unsure),
    player2OpponentPerception: String(human|ai|unsure),
    player1PerceptionConfidence: Number(1-5),
    player2PerceptionConfidence: Number(1-5),
    player1SuspicionTiming: String(round_1_2|round_3_4|round_5_7|round_8_12|round_13_17|round_18_20|never_suspected),
    player2SuspicionTiming: String(round_1_2|round_3_4|round_5_7|round_8_12|round_13_17|round_18_20|never_suspected),
    player1DetectionCues: [String],
    player2DetectionCues: [String],
    player1DetectionOther: String,
    player2DetectionOther: String,
    player1AiAwarenessEffect: String(no_difference|less_persuasive|more_persuasive|not_sure),
    player2AiAwarenessEffect: String(no_difference|less_persuasive|more_persuasive|not_sure),
    player1AiAwarenessJustification: String,
    player2AiAwarenessJustification: String,
    beliefHistory: [{
        round: Number,
        userId: ObjectId,
        player: String(player1|player2),
        belief: String(for|against|unsure),
        beliefValue: Number(0-100),
        influence: Number(0-100),
        confidence: Number(0-100),
        skipped: Boolean,
        createdAt: Date
    }],
    reflections: [{
        round: Number,
        userId: ObjectId,
        paraphrase: String,
        acknowledgement: String,
        timestamp: Date
    }]
  },
  createdAt: Date,
  startedAt: Date,
  completedAt: Date,
  lastActivityAt: Date,
  matchedBy: ObjectId,
  earlyEndVotes: {
    player1Voted: Boolean,
    player2Voted: Boolean,
    humanVoted: Boolean,
    player1Timestamp: Date,
    player2Timestamp: Date,
    expired: Boolean
  },
  completionReason: String
}
\end{verbatim}

\subsubsection{Indexing Strategy}

\begin{table}[htbp]
\centering
\begin{tabular}{lll}
\toprule
\textbf{Collection} & \textbf{Index} & \textbf{Rationale} \\
\midrule
User & username & Guest/user lookup \\
User & isGuest & Filter guest users \\
Debate & player1UserId & User debate history \\
Debate & player2UserId & Opponent debate history \\
Debate & status & Active debate queries \\
Debate & createdAt & Sorting by date \\
\bottomrule
\end{tabular}
\caption{Database Indexes}
\end{table}

\subsection{Real-Time Communication}

\subsubsection{Overview}

Real-time communication is critical for the debate experience and for data fidelity.
Users expect immediate feedback when they submit arguments and instant visibility
when opponents respond, while the study requires low-latency belief capture tied
to specific argument exchanges. Socket.io provides bidirectional communication over WebSocket,
with automatic fallback to polling for incompatible clients.

The communication model is event-based: clients emit events to the server
(e.g., submitArgument), the server processes them and broadcasts updates to
all participants (e.g., argumentAdded). This decouples UI logic from network concerns.

\subsubsection{WebSocket Events}

The platform uses Socket.IO for real-time bidirectional communication.
All events are namespace-scoped to prevent conflicts.
Security is enforced via JWT token validation on connection.

\paragraph{Client to Server Events (Client Emits)}

\begin{enumerate}
    \item \textbf{join:admin}
    \begin{itemize}
        \item Triggered: Admin user connects to dashboard
        \item Payload: None (token passed in handshake auth)
        \item Effect: Validates JWT, verifies admin role, joins socket to \texttt{'admin'} room
        \item Logs: \texttt{[Socket] Client joined admin room: \{socket.id\}}
    \end{itemize}

    \item \textbf{join:debate}
    \begin{itemize}
        \item Triggered: Participant enters active debate session
        \item Payload: \texttt{debateId} (ObjectId string)
        \item Effect: Joins socket to room \texttt{debate:\{debateId\}}, enables debate-specific event delivery
        \item Logs: \texttt{[Socket] Client joined debate room: \{debateId\}}
    \end{itemize}

    \item \textbf{leave:debate}
    \begin{itemize}
        \item Triggered: Participant exits debate or navigates away
        \item Payload: \texttt{debateId} (ObjectId string)
        \item Effect: Removes socket from room \texttt{debate:\{debateId\}}
        \item Logs: \texttt{[Socket] Client left debate room: \{debateId\}}
    \end{itemize}
\end{enumerate}

\paragraph{Server to Client Events (Server Broadcasts)}

\begin{enumerate}
    \item \textbf{debate:created}
    \begin{itemize}
        \item Recipient: \texttt{admin} room
        \item Trigger: New waiting debate created via \texttt{POST /api/debates/join}
        \item Payload: \texttt{\{debateId, gameMode, topicId\}}
        \item Effect: Admin dashboards see new debate in queue
    \end{itemize}

    \item \textbf{debate:started}
    \begin{itemize}
        \item Recipients: Room \texttt{debate:\{debateId\}}, \texttt{admin} room
        \item Trigger: Player2 joins (human-human) or admin matches AI (human-ai)
        \item Payload: \texttt{\{debateId, firstPlayer: 'for' | 'against'\}}
        \item Effect: Both participants notified debate is active; first player announced;
        if AI goes first, opening argument generated
    \end{itemize}

    \item \textbf{debate:argumentAdded}
    \begin{itemize}
        \item Recipients: Room \texttt{debate:\{debateId\}}, \texttt{admin} room
        \item Trigger: Player submits argument via \texttt{POST /api/debates/\{debateId\}/argument}
        \item Payload: \texttt{\{debateId, argument: \{stance, text, round, submittedBy\},
        status, currentRound, nextTurn\}}
        \item Effect: Real-time argument visibility for opponent; turn passes; if opponent
        is AI, response auto-triggered after contextual delay
    \end{itemize}

    \item \textbf{debate:roundAdvanced}
    \begin{itemize}
        \item Recipients: Room \texttt{debate:\{debateId\}}, \texttt{admin} room
        \item Trigger: Both players submit belief updates for current round
        \item Payload: \texttt{\{debateId, currentRound, status, nextTurn\}}
        \item Effect: Round counter increments; if max rounds reached, status
        becomes \texttt{completed}; otherwise next round begins
    \end{itemize}

    \item \textbf{debate:completed}
    \begin{itemize}
        \item Recipients: Room \texttt{debate:\{debateId\}}, \texttt{admin} room
        \item Trigger: Natural completion (20 rounds), mutual early-end agreement,
        admin force-end, or AI vote timeout
        \item Payload: \texttt{\{debateId, reason: String\}}
        \item Effect: UI transitions to post-debate survey; results finalized; admin
        sees debate moved to completed list
    \end{itemize}

    \item \textbf{debate:cancelled}
    \begin{itemize}
        \item Recipients: All connected sockets (global broadcast)
        \item Trigger: Debate creator cancels waiting debate via
        \texttt{DELETE /api/debates/\{debateId\}/cancel}
        \item Payload: \texttt{\{debateId\}}
        \item Effect: Waiting debate removed from all client queues
    \end{itemize}

    \item \textbf{debate:earlyEndVote}
    \begin{itemize}
        \item Recipients: Room \texttt{debate:\{debateId\}}, admin room
        \item Trigger: Player votes to end via \texttt{POST /api/debates/\{debateId\}/vote-end},
        revokes vote, or AI decides to vote
        \item Payload: \texttt{\{debateId, votes: \{player1Voted: Boolean, player2Voted: Boolean\}\}}
        \item Effect: Both players see vote status; if both true in human-human mode,
        debate auto-completes; in human-AI mode, triggers AI handler
    \end{itemize}

    \item \textbf{debates:cleanup}
    \begin{itemize}
        \item Recipients: \texttt{admin} room
        \item Trigger: Scheduled cleanup job detects stale debates (no activity for threshold period)
        \item Payload: \texttt{\{reason: String, cleanedCount: Number\}}
        \item Effect: Admin dashboards refresh to reflect cleaned-up debates
    \end{itemize}

    \item \textbf{Connection/Disconnection Events}
    \begin{itemize}
        \item \texttt{connect}: Socket successfully connected to server
        \item \texttt{disconnect}: Socket disconnected (network, logout, server restart)
        \item \texttt{reconnect}: Socket automatically reconnected after disconnection
        \item \texttt{connect\_error}: Connection failed (token expired, auth denied)
        \item Reconnection: Automatic with exponential backoff (max 5 attempts, 1-5s delay)
    \end{itemize}
\end{enumerate}

\subsubsection{Room Architecture}

\begin{table}[htbp]
\centering
\begin{tabular}{lll}
\toprule
\textbf{Room} & \textbf{Members} & \textbf{Events Received} \\
\midrule
\texttt{admin} & All connected admin sockets & All admin-level events \\
\texttt{debate:\{id\}} & Participants + admin observers & Debate-specific events \\
Global (no room) & N/A & \texttt{debate:cancelled} \\
\bottomrule
\end{tabular}
\caption{WebSocket Room Structure}
\end{table}

\subsubsection{Event Broadcasting Pattern}

Most debate events follow a dual-broadcast pattern for transparency:
\begin{itemize}
    \item \textbf{Debate room broadcast}: Participants in debate receive real-time event
    \item \textbf{Admin room broadcast}: All connected admins receive event for monitoring
    \item \textbf{Exception}: \texttt{debate:created} → \texttt{admin} only;
    \texttt{debate:cancelled} → global broadcast
\end{itemize}

\subsubsection{Event Flow Example: Argument Submission}

\begin{enumerate}
    \item Participant submits argument text via UI
    \item Frontend posts to \texttt{POST /api/debates/\{debateId\}/argument}
    \item Backend validates argument length and turn state
    \item Backend saves argument to debate record in MongoDB
    \item Backend determines next player's turn and stance
    \item Server broadcasts \texttt{debate:argumentAdded} to \texttt{debate:\{debateId\}} room
    \item Opponent's client receives event in real-time, UI updates with new argument
    \item If opponent is AI: Backend schedules AI response generation with contextual delay
    \item When AI response ready: Backend saves and broadcasts another \texttt{debate:argumentAdded}
    \item Round belief check displayed when both arguments submitted
    \item Belief update and reflection captured, timestamped, and persisted to debate record
\end{enumerate}

\subsubsection{Between-Round Belief and Reflection Flow}

\begin{enumerate}
    \item After both arguments in a round, the UI prompts the participant to update stance
    strength, confidence, and perceived influence
    \item Participant provides a brief reflection paraphrasing the opponent's argument
    \item Backend validates input ranges and associates entries with round number and user
    \item Belief and reflection data are stored in \texttt{beliefHistory} and \texttt{reflections}
    \item Completion triggers \texttt{debate:roundAdvanced}, starting the next exchange
\end{enumerate}

\subsubsection{Connection Lifecycle}

\begin{enumerate}
    \item User logs in: Frontend creates socket with JWT token in handshake \texttt{auth}
    \item \texttt{connect} event: Connection successful
    \item Admin user: Emits \texttt{join:admin}; admin room listeners registered
    \item Participant enters debate: Emits \texttt{join:debate} with debateId;
    debate-specific listeners registered
    \item During debate: Receives all events for that debate room
    \item Debate ends: Transitions to post-survey; optionally emits \texttt{leave:debate}
    \item User logs out: Socket disconnects; all rooms automatically cleared
\end{enumerate}

\subsubsection{Transport Configuration}

\begin{itemize}
    \item Primary: \texttt{websocket} (native bidirectional protocol)
    \item Fallback: \texttt{polling} (for environments/proxies blocking WebSocket)
    \item Connection timeout: 20 seconds
    \item Reconnection: Automatic with exponential backoff
    \item EIO3 support: Enabled for legacy client compatibility
\end{itemize}

\subsection{Security Architecture}

Security measures are designed to protect participant anonymity and prevent manipulation of
experimental outcomes while preserving data integrity.

\subsubsection{Authentication}

\begin{itemize}
    \item JWT (JSON Web Tokens) for stateless authentication
    \item Regular user login with username/password via POST /auth/login
    \item Guest user auto-generated usernames via POST /auth/guest-login
    \item Guest session resumption via POST /auth/guest-resume
    \item Token expiration: 24 hours (regular users), 7 days (guests)
    \item No database session records - authentication is purely token-based
\subsection{Turn State Management \& Validation}

Turn validation ensures the protocol is followed and that data are captured at the correct time.
By enforcing turn order and requiring belief checks after each argument pair, the system prevents
missing or misaligned data entries that would compromise round-level persuasion analysis.

\begin{itemize}
    \item Enforces correct alternation of arguments and round boundaries
    \item Blocks progression until belief and reflection data are recorded
    \item Auto-corrects inconsistent state and logs anomalies for auditability
\end{itemize}
    \item User clicks guest login
    \item Backend generates unique username (e.g., "AdorableWolf")
    \item User document created with isGuest=true flag
    \item JWT token generated (7d expiration)
    \item Frontend returns guests list from last 7 days
    \item Returning guest can resume session via POST /auth/guest-resume
\end{enumerate}

\subsection{Scalability Considerations}

\subsubsection{Overview}

The platform's architectural choices prioritize scalability as user load and concurrent debates grow.
Key design decisions—stateless authentication, event-driven real-time updates,
horizontal scaling patterns—ensure we can add more servers without architectural redesign.

This section outlines how the current system can grow, performance bottlenecks to monitor,
and future infrastructure improvements.

\subsubsection{Horizontal Scaling}

\begin{itemize}
    \item Stateless API servers allow load balancing
    \item WebSocket servers require sticky sessions or message queue
    \item Database replication for read scalability
\end{itemize}

\subsection{Error Handling \& Resilience}
    \begin{itemize}
        \item Waiting debates: nextTurn must be null
        \item Active debates (0 args): nextTurn must equal firstPlayer
        \item Active debates (1 arg): nextTurn must be opposite of argument stance
        \item Active debates (2 args): nextTurn must be null (awaiting belief check)
        \item Completed/abandoned debates: nextTurn must be null
    \end{itemize}
\end{itemize}

\textbf{calculateExpectedNextTurn(debate)}
\begin{itemize}
    \item Computes expected nextTurn value from debate state
    \item Uses current round argument count and stances
    \item Returns 'for', 'against', or null
\end{itemize}

\textbf{autoFixTurnState(debate)}
\begin{itemize}
    \item Automatically corrects nextTurn if inconsistent
    \item Modifies debate object in-place
    \item Logs corrections for debugging
    \item Returns true if changes were made
\end{itemize}

\textbf{logTurnStateReport(debate)}
\begin{itemize}
    \item Generates detailed diagnostic report
    \item Logs current vs. expected turn state
    \item Lists validation errors if any
\end{itemize}

\textbf{saveWithValidation(debate, context)}
\begin{itemize}
    \item Safe save wrapper with post-save validation
    \item Logs warnings if validation fails after save
    \item Context parameter for debugging where save occurred
\end{itemize}

\subsubsection{Validation Integration Points}

\paragraph{Debate Model Pre-Save Hook}

The Debate schema includes a pre-save middleware that:
\begin{enumerate}
    \item Updates lastActivityAt timestamp
    \item Validates turn state before database write
    \item Auto-fixes turn state if validation fails
    \item Logs warnings with debate ID and error details
\end{enumerate}

This ensures \textit{all} debate modifications undergo validation, regardless of where
the save originates in the codebase.

\paragraph{AI Response Generation (triggerAIResponse)}

Before generating AI responses:
\begin{enumerate}
    \item Validate current turn state
    \item Auto-fix if inconsistencies detected
    \item Save corrected state before proceeding
    \item Verify it's actually the AI's turn
\end{enumerate}

This prevents the "Turn state mismatch" errors that previously occurred when
automatch failed to initialize nextTurn properly.

\paragraph{Argument Submission}

After human or AI argument submission:
\begin{enumerate}
    \item Debate saved with updated arguments and nextTurn
    \item Post-save validation confirms consistency
    \item Warnings logged if issues detected
\end{enumerate}

\paragraph{AutoMatch AI Assignment}

When automatch assigns AI opponent to waiting debate:
\begin{enumerate}
    \item Update debate status to 'active'
    \item Set firstPlayer randomly ('for' or 'against')
    \item \textbf{Critical:} Initialize nextTurn to firstPlayer
    \item Validate turn state after update
    \item Log warnings if validation fails
\end{enumerate}

This fix resolves the original bug where nextTurn was null after automatch,
causing AI response generation to fail.

\subsubsection{Error Recovery Strategy}

The turn validation system employs a defensive, self-correcting approach:

\begin{enumerate}
    \item \textbf{Detection}: Validation runs at multiple checkpoints (pre-save hook, before AI response, after argument submission)
    \item \textbf{Correction}: auto-fix immediately repairs inconsistent state
    \item \textbf{Logging}: All corrections logged with debate ID and error details for post-mortem analysis
    \item \textbf{Non-blocking}: Validation warnings do not halt operations, allowing system to self-heal
\end{enumerate}

This prevents cascading failures while maintaining audit trail for debugging.

\subsection{Error Handling \& Resilience}

\subsubsection{API Error Handling}

\begin{itemize}
    \item Standardized error response format
    \item HTTP status codes for error classification
    \item Error logging to centralized system
    \item User-friendly error messages in frontend
    \item Turn state validation and auto-correction
\end{itemize}

\subsubsection{Network Resilience}

\begin{itemize}
    \item WebSocket reconnection with exponential backoff
    \item Fallback to polling if WebSocket unavailable
    \item Request retry logic with circuit breaker
    \item Graceful degradation of features
\end{itemize}

\subsubsection{Database Resilience}

\begin{itemize}
    \item Connection pooling and reconnection logic
    \item Automatic retry on transient failures
    \item Transaction support for critical operations
\end{itemize}

\subsection{Operational Considerations for Data Quality}

The study relies on complete, reliable data capture. Operational safeguards focus on
preserving data integrity rather than infrastructure breadth:
\begin{itemize}
    \item Resilient WebSocket reconnection to avoid missing real-time events
    \item Defensive validation and logging for all belief and reflection entries
    \item Rate limits to prevent spam that could contaminate data
    \item Consistent error handling to avoid partial or corrupted debate records
\end{itemize}

\subsection{Conclusion}

\subsubsection{Design Summary}

The Debate Platform architecture successfully balances several competing concerns:

\begin{itemize}
    \item \textbf{Real-time Responsiveness}: WebSocket communication provides instant
    debate updates and low-latency argument exchange
    \item \textbf{Experimental Validity}: Information hiding and structured belief capture
    reduce source bias and enable round-level persuasion measurement
    \item \textbf{Scalability}: Stateless authentication and horizontal scaling
    strategies enable growth without architectural redesign
    \item \textbf{Simplicity}: Event-driven design and modular backend keep code maintainable and testable
    \item \textbf{Flexibility}: MongoDB's document model and JWT tokens reduce rigid schema constraints
    \item \textbf{Robustness}: Multi-layered turn validation with automatic correction
    prevents logic loss and ensures debate flow integrity
    \item \textbf{User Experience}: Both registered and guest users supported seamlessly;
    pre/post-debate surveys capture feedback
\end{itemize}

\subsubsection{Key Achievements}

This architecture successfully integrates:
\begin{itemize}
    \item RESTful API for CRUD operations (debate creation, user queries)
    \item WebSocket for real-time debate interactions (arguments, AI responses)
    \item MongoDB for flexible document storage with indexed queries and analysis-ready records
    \item OpenAI integration for dynamic AI opponent responses
    \item JWT-based stateless authentication enabling horizontal scaling
    \item Automated cleanup jobs for data maintenance
    \item Defensive turn validation system with automatic state correction
    \item Pre-save model hooks preventing database inconsistencies
    \item Enhanced automatch with proper turn state initialization
    \item Between-round belief and reflection capture for fine-grained persuasion analysis
\end{itemize}

The system design supports core functionality while maintaining flexibility for future enhancements,
performance optimizations, and infrastructure scaling. The turn validation architecture demonstrates
defensive programming principles, preventing edge cases from cascading into system failures.